\section{Reasoning 与 verifier：不要过早固定直觉}

访谈发生在 DeepSeek R1 发布后不久。Lample 的评价包含两层：一是肯定公开论文提供成功与失败实验，能节省整个社区的搜索成本；二是提醒 reasoning 仍处早期，今天关于模型大小、RL 方法和 test-time compute 的直觉，可能像 GPT-3 早期直觉一样迅速过时。

\subsection{可扩展 reasoning 的系统单元}

Reasoning training 需要一个环境给出任务和可执行动作，一个 verifier 判断结果，一个采样/搜索策略探索候选，再把 reward 或筛选结果用于更新模型。若没有可靠 verifier，长 chain-of-thought 只能增加文本长度，不能保证正确性。

\lecturefigure{13-reasoning-verifier-loop.png}{Reasoning 的可扩展单元是 environment--policy--verifier--update 闭环。}{本地字幕 27:45--30:45；概念重绘。}

\begin{knowledgebox}{术语消化：reasoning 系统中的四个层次}
\begin{tabularx}{\textwidth}{L{0.18\textwidth}X X}
\toprule
层次 & 核心问题 & 常见失败 \\
\midrule
Environment & 状态、动作和终止条件是否可执行 & 任务不可重置、工具不稳定、状态泄漏 \\
Verifier & 能否判断最终或中间结果 & reward hacking、错误测试、只测格式 \\
Search / sampling & 如何分配 test-time compute & 候选高度相关、成本失控、过早剪枝 \\
Training update & 如何从结果改进 policy & 过拟合 verifier、能力回退、数据偏斜 \\
\bottomrule
\end{tabularx}
\end{knowledgebox}

\begin{warningbox}{公开失败实验也是高价值基础设施}
当论文说明 PRM、MCTS 或某种 reward 设计在其设置中没有工作，它不是“负面新闻”，而是在限定条件下缩小社区搜索空间。前提是记录环境、模型、数据和评测，否则失败无法复现，也不能迁移到别的设置。
\end{warningbox}

\teachervoice{老师没有把 DeepSeek R1 解释成对 Mistral roadmap 的颠覆，而是把开放研究视为可复用资产。更重要的是，他承认 reasoning 时代的直觉仍不稳定。这种“不知道”比事后制造确定性更适合写进讲义。}

\subsection{本章小结}

Reasoning 系统的中心不是更长文本，而是可执行环境和可信 verifier。早期阶段应扩大实验覆盖、公开负结果并保持目标与证据边界，而不是把单一配方神化。

